\documentclass[11pt,a4paper]{article}
\usepackage[T1]{fontenc}
\usepackage[utf8]{inputenc}
\usepackage{lmodern}
\usepackage{amsmath,amssymb}
\usepackage[margin=25mm]{geometry}
\usepackage[hidelinks]{hyperref}
\hypersetup{pdftitle={Time blocking, Newton polygons and explicit quantum action},pdfauthor={navstokgap research working notes}}
\newcommand{\tightlist}{\setlength{\itemsep}{0pt}\setlength{\parskip}{0pt}}
\setlength{\emergencystretch}{3em}
\setcounter{secnumdepth}{0}
\begin{document}
\section{Time blocking, Newton polygons and explicit quantum
action}\label{time-blocking-newton-polygons-and-explicit-quantum-action}

\textbf{First calculation, 2026-10-27; written derivation, unrefereed.}
In one spatial coordinate and one time variable, the same discrete Green
function controls three operations: constructing a Newton polygon on an
arbitrary time partition, solving its nonlinear boundary-value problem,
and bounding the inverse action Hessian used to integrate fluctuations.
A force Lipschitz constant gives bounds uniform in the number and
placement of inserted times. The quadratic fluctuation covariance
contains an explicit factor of \(\hbar\); this first calculation does
not derive its positivity.

\textbf{Targeted second opinion (GPT-6.1 Terra, 2026-10-27).} Terra read
this note and checked the Green row sum, Hessian/response bounds,
Gaussian covariance, oscillator determinant and constant-force signs
without finding an algebraic error. This was a targeted written check,
not a full independent proof review. The arbitrary-cut interpretation
below came from that pass.

The user directs Newton action necessity first, beginning with Balaban's
techniques in \(0+1\) dimensions. The comparison with Alpoge's work and
Yang--Mills is deferred. Here \(0+1\) means a mechanical degree of
freedom along time, not a field with spatial lattice sites. The
construction is a modern force polygon; no identification with every
historical Newton update is asserted. The Galileo observable remains the
inertial-line / falling-parabola defect, not swept sectors.

We assemble standard Green-function, contraction and
Gaussian-elimination arguments. No literature novelty is claimed. The
source ingredient is Balaban's controlled effective density under
blocking, described in his
\href{../docs/balaban/Balaban_UltravioletStability_ICM1986_1987.md}{ICM
overview}, printed pp.~1260--1262 (text read), and Dimock's
\href{../docs/balaban/Dimock_BalabanI_1108.1335v2.md}{small-field
exposition}, sections 2.1, 2.4, 4.5 and Theorem 14 (selected passages
read). The \href{refinement-composition-and-limit.md}{existing
refinement note} constructs constant-force and harmonic limits. The
\href{cut-measure-newton.md}{cut-measure note}, Proposition 7, already
uses nodal exactness of the same one-dimensional Green function. The
\href{rivero-1998-conjecture-central-forces.md}{Rivero test}, section 2,
supplies the determinant/Schur identity and its branch caveat.

\subsection{1. Actions and the constant to be
investigated}\label{actions-and-the-constant-to-be-investigated}

Let \(m,T>0\), \(x,y\in\mathbb R\), and \(\pi=(0=t_0<\cdots<t_N=T)\).
Write \(\tau_j=t_{j+1}-t_j\) and \(w_j=(\tau_{j-1}+\tau_j)/2\) for
interior vertices. For a potential \(V\) define

\[
 S_\pi^\sigma(q)=\sum_{j=0}^{N-1}
 \frac{m(q_{j+1}-q_j)^2}{2\tau_j}
 +\sigma\sum_{j=0}^{N-1}\frac{\tau_j}{2}
       [V(q_j)+V(q_{j+1})],\qquad q_0=x,\ q_N=y.
 \tag{1}
\]

The real-time Newton action is \(S_\pi^R=S_\pi^{-1}\), with weight
\(e^{iS_\pi^R/\hbar}\) when \(\hbar>0\). The Euclidean action is
\(S_\pi^E=S_\pi^{+1}\), with weight \(e^{-S_\pi^E/\hbar}\). Their
stationary equations have opposite potential signs. They must not be
interchanged when interpreting a background as a Newton trajectory. At
\(\hbar=0\) these exponential expressions are not defined; the classical
stationary equation is defined directly. No relativistic premise or
\(c\) is used in this calculation. Ordinary Planck action is
\(h_P=2\pi\hbar\).

For real time put \(f=-V'/m\), the acceleration. Section 2 needs only a
Lipschitz function \(f\), even without first specifying \(V\). A
Lipschitz condition is a sufficient uniqueness and stability condition;
continuity alone can give local existence. Neither statement alone
guarantees continuation for all times or uniqueness of a two-endpoint
problem.

\subsection{2. Arbitrary intermediate
times}\label{arbitrary-intermediate-times}

Let \(\ell(t)=x+(y-x)t/T\) and

\[
 G(t,s)=\min(t,s)-\frac{ts}{T}.
 \tag{2}
\]

This is the positive Dirichlet Green function of \(-d^2/dt^2\). Define
the polygon map, with piecewise-linear output, by

\[
 (\mathcal T_\pi q)(t)=\ell(t)
       -\sum_{j=1}^{N-1}w_jG(t,t_j)f(q(t_j)).
 \tag{3}
\]

\textbf{Lemma 1 (the exact row sum).} At every vertex,

\[
 \sum_{j=1}^{N-1}w_jG(t_i,t_j)=\frac{t_i(T-t_i)}2.
 \tag{4}
\]

At every \(t\in[0,T]\) the sum is at most \(T^2/8\).

\emph{Proof.} For fixed \(t_i\), the function \(s\mapsto G(t_i,s)\) is
linear on each partition interval and vanishes at the endpoints. Its
trapezoidal integral is therefore exact, and its integral is
\(t_i(T-t_i)/2\). As a function of \(t\), the sum in (4) is the linear
interpolant of that concave quadratic. It lies below the quadratic,
whose maximum is \(T^2/8\). This proof uses no ratio between consecutive
time steps. \(\square\)

Assume \(|f(u)-f(v)|\le\Lambda|u-v|\) for all \(u,v\), and put

\[\theta=\frac{\Lambda T^2}{8}<1.\tag{5}\]

\textbf{Proposition 2 (one bound for every polygon).} The map (3) has a
unique fixed polygon. Its slopes \(v_j=(q_{j+1}-q_j)/\tau_j\) satisfy

\[v_j-v_{j-1}=w_jf(q_j).\tag{6}\]

The continuous map

\[
 (\mathcal Tq)(t)=\ell(t)-\int_0^T G(t,s)f(q(s))\,ds
 \tag{7}
\]

has a unique fixed point, which is \(C^2\) and solves \(\ddot q=f(q)\)
with the two prescribed endpoints. Every sequence of partitions with
\(|\pi|\to0\) converges to this solution. Neither nesting nor equal
intervals is required.

\emph{Proof.} Lemma 1 bounds the Lipschitz constants of (3) and (7) in
the supremum norm by \(\theta\). Repeated application converges because
the successive differences form a geometric series. This proves
existence and uniqueness in the complete polygon space or in
\(C([0,T])\). The derivative of \(G(\cdot,t_j)\) jumps by \(-1\) at
\(t_j\), giving (6). Differentiating (7) gives the continuous equation.

Inserting a vertex changes its neighbours' weights \(w_j\) as well as
the new solution. It does not add one kick while freezing the old force
history. Retained vertices need not remain fixed when the trapezoidal
action itself is refined.

For a quantitative limit, write \(B=\|f(q)\|_\infty\) and
\(U=\|\dot q\|_\infty\) for the continuous solution. For every fixed
\(t\), \(s\mapsto G(t,s)f(q(s))\) is Lipschitz with constant at most
\(B+T\Lambda U/4\). On an interval of length \(\tau\), replacing its
integral by the endpoint trapezoid has error at most half that constant
times \(\tau^2\). Hence

\[
 \|q_\pi-q\|_\infty\le
 \frac{T|\pi|}{2(1-\theta)}
       \left(B+\frac{T\Lambda U}{4}\right).
 \tag{8}
\]

Indeed, subtract the two fixed-point equations and absorb the
\(\theta\|q_\pi-q\|_\infty\) term. The constants in (8) belong to the
fixed solution and interval, not to the partition. \(\square\)

For a force Lipschitz only in a tube \(|q(t)-\ell(t)|\le r\), the same
proof holds within that tube if \(\theta<1\) and

\[
 \frac{T^2}{8}\big(B_0+\Lambda r\big)\le r,
 \qquad B_0=\sup_t|f(\ell(t))|.
 \tag{9}
\]

This makes the closed tube invariant. Uniqueness is then asserted only
inside it. Equation (9) is a finite version of a small-field domain that
survives the construction; it cannot justify discarding paths outside it
from a statistical or oscillatory integral.

\subsection{3. The action Hessian and
elimination}\label{the-action-hessian-and-elimination}

Let \(\mathsf K_\pi\) be the Dirichlet stiffness matrix, with diagonal
\(1/\tau_{j-1}+1/\tau_j\) and adjacent entries \(-1/\tau_j\), and let
\(\mathsf W_\pi=\operatorname{diag}(w_j)\). The tent function (2) gives

\[
 (\mathsf K_\pi^{-1})_{ij}=G(t_i,t_j),\qquad
 \|\mathsf K_\pi^{-1}\mathsf W_\pi\|_{\infty\to\infty}
       \le\frac{T^2}{8}.
 \tag{10}
\]

The first identity follows by taking the slope jump of that tent at each
node; the second is Lemma 1. In particular, for \(V\in C^2\) with
\(|V''|\le m\Lambda\), the Hessian of (1) obeys

\[
 \mathsf H_\pi^\sigma
   =m\mathsf K_\pi+\sigma\mathsf W_\pi
             \operatorname{diag}(V''(q_j))
   \ \ge\ m(1-\theta)\mathsf K_\pi.
 \tag{11}
\]

To see the quadratic-form inequality, the nonnegative matrix
\(\mathsf K_\pi^{-1}\mathsf W_\pi\) has spectral radius at most its
maximum row sum. It is similar to the positive symmetric matrix
\(\mathsf W_\pi^{1/2}\mathsf K_\pi^{-1}\mathsf W_\pi^{1/2}\). Thus
\(\mathsf W_\pi\le(T^2/8)\mathsf K_\pi\), which proves (11). The factor
\(1-\theta\) is uniform over arbitrary partitions.

For a physical force source \(J_j\), inserted as a linear action term
\(\sum_jw_jJ_jq_j\), the stationary response therefore has the bound

\[
 \|\delta q\|_\infty\le
 \frac{T^2}{8m(1-\theta)}\|\delta J\|_\infty.
 \tag{12}
\]

Here the bound is for the differential response. It follows directly by
factoring \(\mathsf H_\pi^\sigma\) as
\(m\mathsf K_\pi(I+\mathsf K_\pi^{-1}\mathsf W_\pi
\sigma\operatorname{diag}(V'')/m)\) and using a Neumann series.

With the global curvature bound, (11) makes the finite action coercive
and strictly convex in its interior vertices. Eliminating some vertices
by minimization is then unambiguous, and its Hessian is the Schur
complement

\[
 \mathsf H_{BB}-\mathsf H_{BA}\mathsf H_{AA}^{-1}\mathsf H_{AB}.
 \tag{13}
\]

Successive eliminations associate. With only a local tube, this is a
statement on the chosen nondegenerate branch. In particular, for the
real-time positive quartic potential the unrestricted action is
unbounded below, as the Rivero note already records; local control is
not a global minimization theorem. Nor does the eliminated action
generally retain only the original kinetic and potential terms.

\subsection{\texorpdfstring{4. What the same inverse says when \(\hbar\)
is
present}{4. What the same inverse says when \textbackslash hbar is present}}\label{what-the-same-inverse-says-when-hbar-is-present}

At a Euclidean background, its Gaussian fluctuation reference has
density proportional to \(\exp[-\eta^T\mathsf H_\pi^E\eta/(2\hbar)]\)
and covariance

\[
 \mathsf C_{\pi,\hbar}^E=\hbar(\mathsf H_\pi^E)^{-1},\qquad
 (\mathsf C_{\pi,\hbar}^E)_{jj}
 \le\frac{\hbar}{m(1-\theta)}\frac{t_j(T-t_j)}{T}
 \le\frac{\hbar T}{4m(1-\theta)}.
 \tag{14}
\]

These are exact Gaussian statements. For a nonlinear action they do not
assert the covariance bound for its full measure without a further
argument. In real time, normalized quadratic Fresnel integration gives
the algebraic centered second moment \(i\hbar(\mathsf H_\pi^R)^{-1}\),
which is not a probability covariance.

For one block of Gaussian variables, integration retains both the Schur
complement (13) and a determinant factor: respectively
\((2\pi\hbar)^{|A|/2}(\det\mathsf H_{AA}^E)^{-1/2}\) in Euclidean time
and \((2\pi i\hbar)^{|A|/2}(\det\mathsf H_{AA}^R)^{-1/2}\) when the real
Hessian is positive. Determinants, normalization and any higher
interactions must accompany the stationary action. These are the extra
data that Balaban's effective-density viewpoint instructs us to retain.

Equations (11)--(14) identify the connection to ODE control: the
curvature constant that makes the background construction stable also
controls the Gaussian fluctuations for every positive supplied
\(\hbar\). They do not create the multiplier \(\hbar\) in (14).

\subsection{5. Written controls: Galileo and the
oscillator}\label{written-controls-galileo-and-the-oscillator}

\textbf{Constant force.} For \(V(q)=-Fq\), \(\Lambda=0\). Equation (7)
is exactly

\[q(t)=\ell(t)+\frac{F}{2m}t(t-T).\tag{15}\]

By (4), every polygon has these same vertex positions. The fluctuation
Hessian is the free one for either action sign, and the Euclidean
covariance is \((\hbar/m)G(t_i,t_j)\). For a single inserted midpoint of
a cell of length \(T\), its variance is \(\hbar T/(4m)\); for an unequal
cut at \(u\) it is \(\hbar u(T-u)/(mT)\). The deterministic vertex
survives at \(\hbar=0\); the Gaussian width tends to zero. This uses
precisely the Galileo comparison already normalized in the refinement
note.

\textbf{One arbitrary cut (Terra's connection).} Split \(T=u+v\) and let
\(H_u\) be the piecewise-linear hat of height one at \(u\), zero at the
cell endpoints. Then

\[
 \int_0^T(H_u')^2dt=\frac{T}{uv},\qquad
 a_F=-\frac{Fuv}{2m},\qquad
 \frac m2\int_0^T(a_FH_u')^2dt
       =\frac{F^2uvT}{8m}=K_T-K_u-K_v.
 \tag{15a}
\]

where \(K_s=F^2s^3/(24m)\). The free bridge's conditional
inserted-vertex innovation has variance \(s_u^2=\hbar uv/(mT)\). For the
two equal-variance Gaussian vertex laws, with means separated by
\(a_F\), direct integration of their log-density ratio gives

\[
 D\big(N(a_F,s_u^2)\,\|\,N(0,s_u^2)\big)
       =\frac{a_F^2}{2s_u^2}
       =\frac{K_T-K_u-K_v}{\hbar}.
 \tag{15b}
\]

Here \(D\) is relative entropy in nats. Successive bridge innovations
factor into independent Gaussian hat coefficients: conditioned on a
cell's endpoints, its inserted value has the variance above, and the two
residual bridges are independent. The hats are also orthogonal in the
Dirichlet energy because each finer hat has zero derivative integral
inside a cell where an older hat has constant derivative.

Thus classical vertex elimination, conditional Gaussian normalization
and the cut energy are related by one Schur factorization. The
cut-measure note already proves the additive action budget; the useful
connection here is to the blocking variables. Equation (15b) assumes
these Gaussian laws. It does not establish a physical bound on the
information in complete apparatus records.

\textbf{Oscillator and caustics.} Let \(V=m\omega^2q^2/2\),
\(\omega>0\). Its initial-value equation has the unique explicit
solution \(q(t)=q(0)\cos\omega t+\dot q(0)\sin\omega t/\omega\) for all
times. Nevertheless its two-endpoint problem loses uniqueness at
\(T=\pi/\omega\) when \(y=-x\); for other \(y\) it has no solution at
that time. A singular boundary Hessian therefore need not mean a failure
of initial-value mechanics.

On a uniform \(N\)-cell partition, \(\tau=T/N\), \(N\ge2\), the first
Dirichlet eigenvalue of the difference operator is
\(4\tau^{-2}\sin^2(\pi/(2N))\). Its Hessian first becomes singular at

\[
 T_N=\frac{2N}{\omega}\sin\frac{\pi}{2N}
       \longrightarrow\frac\pi\omega.
 \tag{16}
\]

This checks the partition-independent sufficient window
\(\omega^2T^2<8\): \(T_2=\sqrt8/\omega\), while the continuum caustic is
later. An ODE Lipschitz constant cannot exclude those later caustics.

For an exact Gaussian control, suppose \(T<\sqrt8/\omega\) and set
\(\alpha=2\arcsin(\omega\tau/2)\). The determinant ratio is

\[
 R_N=\frac{\det\mathsf H_\pi^R}{\det(m\mathsf K_\pi)}
     =\frac{\sin(N\alpha)}{N\sin\alpha}
       \longrightarrow\frac{\sin(\omega T)}{\omega T}.
 \tag{17}
\]

Indeed the determinant of the tridiagonal matrix with diagonal
\(2-\omega^2\tau^2=2\cos\alpha\) and off-diagonal \(-1\) satisfies
\(D_n=2\cos\alpha D_{n-1}-D_{n-2}\), with \(D_0=1,D_1=2\cos\alpha\).
Thus \(D_{N-1}=\sin(N\alpha)/\sin\alpha\); the free determinant is
\(N\). For \(N=2\), (17) gives \(R_2=1-\omega^2\tau^2/2\), also obtained
by completing the single midpoint square.

With the usual free normalization for each cell, finite Gaussian
integration gives the real-time kernel

\[
 K^\pi_\hbar(y,x)=
 \sqrt{\frac{m}{2\pi i\hbar T}}\,R_N^{-1/2}
           e^{iS_\pi^R(q_\pi)/\hbar}.
 \tag{18}
\]

The prefactor follows from \(\det\mathsf K_\pi=T/\prod_j\tau_j\) and
completing the quadratic form. The determinant identity follows by the
same tridiagonal recurrence, or by successive free one-vertex
eliminations. Equation (18) retains the same \(\hbar\) at every depth.
Reversing the sign of the oscillator term would put
\(1+\omega^2\tau^2/2\) in the two-cell determinant and remove the
real-time conjugate point. This sign test distinguishes the Euclidean
and Newton constructions.

\subsection{6. Large fields, continuation and the next scale
question}\label{large-fields-continuation-and-the-next-scale-question}

A local Lipschitz bound controls a local construction. To continue an
initial-value solution, one also needs control preventing escape from
every bounded phase-space region. In the confining example
\(V(q)=\lambda q^4/4\), \(\lambda>0\), conservation of a fixed classical
energy \(E>0\) gives

\[
 |q|\le R_E=(4E/\lambda)^{1/4},\qquad
 |p|\le\sqrt{2mE},\qquad
 |f'|\le\frac{3\lambda R_E^2}{m}.
 \tag{19}
\]

These bounds permit continuation of the initial-value solution. Applying
the local polygon lemma to a whole tube may require a larger radius than
\(R_E\) and the invariance condition (9). For the inverted quartic,
\(q(t)=q_0/[1-q_0\sqrt{\lambda/(2m)}t]\), \(q_0>0\), is instead a
finite-time escape solution. Its force is locally Lipschitz everywhere.
This is an explicit reason to distinguish local control from
continuation. Balaban's large-field estimates address weighted integrals
over all field regions; that is a statistical obligation distinct from
trajectory continuation.

Keeping \(\hbar\) explicit in the positive quartic Euclidean action, set
\(t=Ts\) and \(q=\sqrt{\hbar T/m}\,z\). Then

\[
 \frac{S^E}{\hbar}
 =\int_0^1\left[\frac{(z')^2}{2}+\frac g4z^4\right]ds,
 \qquad g=\frac{\lambda\hbar T^3}{m^2}.
 \tag{20}
\]

For the zero-endpoint free reference, the Brownian-bridge reflection
bound is

\[
 \mathbb P\!\left(\sup|q|>b\sqrt{\hbar T/m}\right)
       \le2e^{-2b^2},\qquad b>0.
 \tag{21}
\]

It follows by applying the positive and negative one-sided reflection
bounds and a union bound; the Rivero note, section 3, supplies the same
calculation. On this small-field tube the real quartic curvature gives
\(\theta\le3b^2g/8\). Thus the natural perturbation parameter combines
force curvature, the observation interval and the explicitly supplied
quantum width. For the interacting normalized measure, its denominator
and its own large-field contribution still require control; (21) alone
is not that result.

The next useful test is to retain the normalized fluctuation and
large-field terms through repeated time insertion, and determine which
physical premise would fix their action multiplier independently of the
weight \(e^{-S^E/\hbar}\) or \(e^{iS^R/\hbar}\). The existing
\href{rivero-1998-conjecture-central-forces.md}{quartic zero-limit
theorem} prevents treating refinement consistency alone as that premise.
A loss of boundary invertibility also needs the oscillator countertest
above before it is used to argue for completion of Newton's dynamics.

A complementary construction suggested by Terra is to control the
initial-value polygon flow and its tangent map under arbitrary
insertion. Their composition is causal and remains meaningful through
the oscillator's boundary caustics. Comparing that description with the
endpoint action would separate loss of one coordinate description from
an actual need for additional physical dynamics.

\subsection{7. Consequence for STATE}\label{consequence-for-state}

Newton action necessity is the active priority. This note gives the
first \(0+1\) background construction: arbitrary unequal polygons
converge on a stated Lipschitz-controlled interval; the same Green
function controls source response and the Gaussian width proportional to
\(\hbar\). Its oscillator example separates boundary caustics from
initial-value failure. These are written, unrefereed calculations. Next
retain the full weighted corrections and justify their physical
weighting without importing the positive action constant to be derived.
Alpoge and Yang--Mills remain for a later comparison.
\end{document}
